% =========================
% CHAPTER 2 — BACKGROUND & FUNDAMENTALS (CLEAN + WITH EXTRA (NEW) CITATIONS)
% File: chapter2.tex
% =========================

\chapter{Background and Fundamentals} \label{chap:ch2}

This chapter introduces the minimum set of oceanographic and robotic fundamentals required to
frame the research problem addressed in this report. The emphasis is on concepts (not on specific
solutions from the literature): (i) ocean features that naturally induce distinct sampling regimes,
(ii) AUV operational constraints that shape feasibility, (iii) a high-level problem formulation for
closed-loop adaptive sampling, (iv) common objective functions used to quantify sampling utility,
and (v) essential machine-learning notions for regime inference. Detailed comparisons of published
methods are deferred to Chapter~\ref{chap:ch3}. The discussion is aligned with the operational
oceanography view of coupling observations, modelling, and forecasting skill \parencite{guo2025orcadl,balaji2022gcmobsolete}.

\section{Ocean features relevant to sampling} \label{sec:ocean_features}

Coastal and shelf seas exhibit strong spatial and temporal heterogeneity. For adaptive sampling,
heterogeneity matters because it creates regions with different dynamics, predictability, and
observational signatures, motivating the use of \emph{regimes} as an organising concept. This
heterogeneity is also central in coastal ecological and biogeochemical contexts, where patchiness
and gradients shape primary production and ecosystem processes \parencite{falkowski1998biogeochemical,martin2003patchiness}.
Typical AUV sensor payloads used to observe such regimes are summarised in Table~\ref{tab:sensors_phenomena}.

\subsection{Water masses and vertical structure (thermocline/halocline)}
A common description of ocean structure uses \emph{water masses}: bodies of water with coherent
physical properties. In practice, water masses are often distinguished by their temperature--salinity
(T--S) characteristics and by vertical structure. A sharp vertical temperature gradient defines a
\emph{thermocline}, while a sharp salinity gradient defines a \emph{halocline}. These layers can
shift in depth and intensity, influencing mixing and transport.

For sampling, vertical structure matters because (i) the same horizontal location may correspond to
different regimes depending on depth, and (ii) strong stratification can reduce coupling across layers,
changing how informative a profile is for updating a model.

\subsection{Fronts and boundary structures}
Ocean \emph{fronts} are narrow transition zones separating water masses with different properties,
typically expressed as strong horizontal gradients in temperature and salinity. Fronts are key
sampling targets because they can evolve rapidly and often represent informative boundaries whose
mischaracterisation can propagate into downstream modelling and decision-making. Feature- and
boundary-driven robotic sampling in coastal waters has been explicitly motivated in prior adaptive
ocean sampling work \parencite{fossum2019phytoplankton}.

\subsection{Current-driven structures as regime indicators}
Beyond T/S signatures, ocean regimes may also be \emph{indicated} by the flow field. Currents organise transport and
mixing into coherent structures (e.g., jets, eddies, and filaments) that shape gradients and tracer accumulation,
thereby partitioning the domain into dynamical zones \parencite{weiss2007transportmixing,mcgillicuddy2016mechanisms,mahadevan2016submesoscale}.
When velocity information is available (from forecasts or onboard estimates), this flow context provides additional
intuition for regime boundaries and motivates flow-aware feasibility considerations in time-varying environments
\parencite{aguiar2018trajectoryopt}. Specific diagnostics and their use in regime identification are reviewed in
Chapter~\ref{chap:ch3}.






\subsection{Plumes and river-influenced coastal dynamics}
River plumes introduce buoyant, fresher water that interacts with tides, wind forcing, and shelf
circulation, creating sharp salinity gradients and distinct optical/biogeochemical signatures.
Plumes can shift direction and extent on hourly-to-daily time scales, providing a clear example of a
heterogeneous environment where sampling across plume and ambient waters may encounter distinct regimes.
Such plume- and bloom-influenced settings are also a canonical motivation for adaptive robotic sampling
in coastal observing systems \parencite{fossum2019phytoplankton}.

\subsection{Temporal variability, uncertainty, and why both matter}
Temporal variability arises from tides, wind-driven circulation, diurnal heating/cooling, and episodic
events. A practical implication is that sampling value is time-dependent: a location can be informative
now but not later. Two concepts are useful:

\begin{itemize}
  \item \textbf{Uncertainty}: what the model does not know (e.g., forecast variance/spread), often addressed through
  data assimilation and related update mechanisms \parencite{dowd2014dataassimilation}.
  \item \textbf{Variability}: how fast the environment changes (e.g., moving boundaries), which in patchy coastal
  systems can dominate where/when observations matter most \parencite{martin2003patchiness}.
\end{itemize}

Adaptive strategies that focus only on uncertainty may under-sample fast-changing boundaries,
while strategies that ignore uncertainty may waste effort where the model is already confident.
Chapter~\ref{chap:ch3} discusses how different works operationalise these ideas.

\subsection{Why temperature, salinity, and depth are strong ``axes'' for regimes}
Temperature and salinity jointly determine density and stratification, which shape circulation and mixing.
Together with depth, they provide a physically meaningful coordinate system for regimes:
(i) T/S profiles encode water-mass identity; (ii) depth indexes layered structure; and (iii) gradients
in T/S highlight fronts and boundaries. These variables also align well with typical AUV sensing (CTD payloads).

\section{AUV sampling and constraints} \label{sec:auv_constraints}

Adaptive sampling must operate under real AUV constraints. These constraints shape both feasibility
and robustness of planning policies. A general overview of constraints and their implications for
adaptive sampling design is provided in survey work on AUV adaptive sampling methods
\parencite{hwang2019auvreview}. At a systems level, the broader push toward autonomous ocean sampling
networks further emphasises robustness, persistence, and operational integration \parencite{curtin2009aosn}.

\subsection{Platform and mission constraints}
\textbf{Energy and endurance:} Finite energy budgets constrain mission duration, range, and sampling intensity.
Energy use depends on speed, manoeuvres, payload operation, and currents.

\textbf{Time and scheduling:} Missions are often deadline-driven (e.g., data must be available before a model update),
which couples routing and scheduling.

\textbf{Communication limitations:} Underwater communication is intermittent and low throughput, so continuous coordination
cannot be assumed; plans must tolerate delays and missed messages.

\textbf{Safety and navigability:} Obstacles, bathymetry, restricted areas, and safety margins impose geofences and may make
transitions infeasible.

\textbf{Time-varying currents:} Currents affect ground-relative motion, travel time, and energy expenditure; in time-varying
flows, the cost of moving between locations depends on departure time, motivating flow-aware planning
\parencite{aguiar2018trajectoryopt}.

\subsection{Multi-agent sampling: centralised vs distributed coordination}
With multiple AUVs, coordination is required to avoid redundancy and satisfy fleet-wide constraints.
Two broad paradigms are commonly used:

\begin{itemize}
  \item \textbf{Centralised}: a coordinator assigns tasks/routes using global information when connectivity permits.
  \item \textbf{Distributed}: each AUV plans locally with limited coordination messages, improving robustness to sparse communications.
\end{itemize}

Hybrid schemes are common in practice (periodic coordination plus local replanning). From an operational perspective,
organisational and architectural choices (e.g., who coordinates what, and when) are increasingly treated as first-class
design decisions in marine autonomy \parencite{skaugset2025autonomousorg}.

\section{High-level problem formulation} \label{sec:problem_formulation}

This section introduces a lightweight formulation that captures the essence of closed-loop, regime-aware,
multi-agent adaptive sampling without committing to a specific estimator or planner.

\subsection{System state and environment}
Consider a team of $N$ AUVs indexed by $i \in \{1,\dots,N\}$. At time $t$, the mission-level state of AUV $i$ is
\begin{equation}
  s_i(t) = \big(p_i(t),\, \theta_i(t),\, E_i(t),\, \tau_i(t)\big),
\end{equation}
where $p_i(t)$ is position (2D or 3D), $\theta_i(t)$ denotes pose-related variables, $E_i(t)$ is remaining energy,
and $\tau_i(t)$ encodes time information (e.g., elapsed time or time-to-deadline).

The ocean field of interest is modelled as a spatiotemporal process
\begin{equation}
  x(p,t),
\end{equation}
where $p$ denotes spatial location (including depth if needed) and $x$ may represent temperature, salinity,
or derived quantities. A forecasting system typically provides a prior estimate $\hat{x}(p,t)$ and a representation
of uncertainty (e.g., variance maps or ensemble spread) \parencite{guo2025orcadl}.

\subsection{Measurements and uncertainty}
When an AUV samples at location $p$ and time $t$, it obtains an observation
\begin{equation}
  y = x(p,t) + \varepsilon,
\end{equation}
where $\varepsilon$ represents measurement noise and modelling error. Observations can be point samples,
vertical profiles, or short trajectory segments.

In heterogeneous environments, the mapping from observations to model updates may depend on regime:
conditioning inference and updates on regime information can help avoid undesired coupling across boundaries.
How published works address this is reviewed in Chapter~\ref{chap:ch3}.

\subsection{Mission-level objective}
At a high level, the mission chooses actions $a_i(t)$ (e.g., waypoints, sampling depths, patterns) to maximise
expected utility over a horizon:
\begin{equation}
  \max_{\{a_i\}} \; \mathbb{E}\Big[ U(\text{information gained}, \text{forecast improvement}) - C(\text{energy}, \text{time}, \text{risk}) \Big],
\end{equation}
subject to feasibility constraints (energy bounds, safety constraints, communication opportunities, and
reachability influenced by currents). In practice, decision support is often framed explicitly in terms
of a value/utility function that trades off scientific gain against operational cost \parencite{aguiar2020valuefunction}.

\section{Fundamentals of adaptive sampling objectives} \label{sec:adaptive_objectives}

Adaptive sampling differs largely in how it quantifies sampling value. Common objective families include:

\subsection{Uncertainty reduction}
If the field estimate admits a probabilistic representation, actions can be scored by expected reduction of posterior
uncertainty (often approximated via variance maps):
\begin{equation}
  U_{\text{var}}(p) \approx \sigma^2_{\text{prior}}(p) - \mathbb{E}\big[\sigma^2_{\text{post}}(p)\big].
\end{equation}
The value of moving beyond static grid surveys toward uncertainty-aware designs has been emphasised in optimal
synoptic survey analyses \parencite{frolov2014grid}.

\subsection{Information gain (entropy / mutual information)}
Information-theoretic objectives quantify expected entropy reduction. For random variables $X$ (field state) and
$Y$ (future observation), mutual information
\begin{equation}
  I(X;Y) = H(X) - H(X \mid Y)
\end{equation}
captures expected uncertainty reduction while accounting for correlation, which helps avoid redundant sampling.

\subsection{Forecast improvement (model-aware objectives)}
When a forecasting model and an update mechanism exist, sampling utility can be defined by expected improvement
in forecast skill (e.g., reduced prediction error on a future horizon), rather than by instantaneous uncertainty reduction.
In biogeochemical and physical forecasting contexts, this objective is tightly linked to assimilation/update methodology
\parencite{dowd2014dataassimilation}.

\subsection{Event/regime-oriented objectives}
For missions focused on detecting or tracking fronts, plumes, or regime boundaries, utility can be defined via
classification confidence or boundary uncertainty, motivating coupling between regime inference and planning.
This viewpoint is particularly natural in coastal ecological monitoring where targets are patchy, dynamic features
\parencite{fossum2019phytoplankton}.

\section{Machine-learning basics for regime detection} \label{sec:ml_regime_basics}

This section states minimal ML concepts used later for the literature review.

\subsection{Clustering, classification, and segmentation}
\textbf{Clustering} groups observations into clusters without labels. \textbf{Classification} assigns observations to
predefined classes (often probabilistically). \textbf{Segmentation} assigns labels over a spatial domain (e.g., a grid),
revealing boundaries between regimes.

\subsection{Online vs offline inference}
\textbf{Offline} methods fit/train on a fixed dataset before deployment. \textbf{Online} methods update incrementally as
new data arrive, enabling regime beliefs to evolve during execution and directly influence replanning.

\subsection{Compact representations}
Regime inference can benefit from compact representations of measurements:
handcrafted features (e.g., thermocline depth), learned embeddings (e.g., encoder outputs), or sparse representations.
Chapter~\ref{chap:ch3} discusses how the reviewed literature uses compact models for regime/scene structure.

\subsection{Regime inference as a probabilistic belief}
Many regime models output a belief over regimes rather than a hard label. Such beliefs are useful for planning:
vehicles can focus sampling near uncertain boundaries, and updates can be conditioned on regime probabilities.
Sequential decision-making under uncertainty is also commonly formalised through stochastic models that update
as new observations arrive \parencite{berget2023dynamicstochastic}.

\begin{table}[t]
  \centering
  \caption{Examples of common AUV-relevant sensors/variables and the ocean phenomena they help capture.}
  \label{tab:sensors_phenomena}
  \begin{tabular}{p{0.22\linewidth} p{0.30\linewidth} p{0.40\linewidth}}
    \hline
    \textbf{Sensor / variable} & \textbf{Typical output} & \textbf{Phenomena captured (examples)} \\
    \hline
    CTD & $T(z)$, $S(z)$, density/stratification & Water masses; thermocline/halocline; frontal gradients; plume freshwater signatures \\
    Fluorometer & Chlorophyll-a fluorescence & Productivity/blooms; plume biogeochemical signatures \\
    Turbidity / optical backscatter & Suspended particles & River plumes; sediment fronts; resuspension events \\
    Oxygen optode & Dissolved oxygen & Ventilation/mixing; hypoxia; provenance indicators \\
    ADCP / DVL current estimates & Current profiles / velocity & Currents and shear; flow constraints for trajectory planning \\
    Nitrate / pH / CDOM (if available) & Biogeochemical tracers & Biogeochemical regime separation; plume influence \\
    \hline
  \end{tabular}
\end{table}

\section{Summary} \label{sec:ch2_summary}

This chapter provided the foundational concepts needed to interpret the state-of-the-art review.
It introduced regime-relevant ocean features (water masses, stratification, fronts, plumes), key AUV
constraints (energy, time, communications, safety, and time-varying currents), a high-level problem
formulation for closed-loop adaptive sampling, common objective families, and minimal ML notions
for regime inference. Chapter~\ref{chap:ch3} builds on this baseline to critically review how published
work instantiates these building blocks and where gaps remain.
